\section{Software Requirements Specification}
\label{ch:srs}

\subsection{Purpose and Scope of the SRS}
This chapter condenses the IEEE 830 Software Requirements Specification delivered in Phase~II (\texttt{docs/SRS\_EcoMesh\_Group20.pdf}). It states \emph{what} EcoMesh must do, not how. The scope is a single deployment of cooperating nodes with ten system features (peer discovery, message composition, multi-hop forwarding, adaptive routing, priority handling, reliability, store-and-forward, message security, logging, rescue dashboard); no server exists in any operating mode.

\subsection{Intended Users and Stakeholders}
The \emph{trapped user} originates an SOS or a message and needs one prominent control and no configuration. The \emph{relay user} passively carries other people's traffic and must not be able to read it. The \emph{gateway operator} holds the residual outside link. The \emph{rescue coordinator} monitors the SOS queue, the topology and delivery statistics. The \emph{administrator} assigns node identifiers and constrains who hears whom. The project guide and evaluators review scope, rigour and evidence.

\subsection{Overall Product Description}
Every node runs the same program. HELLO beacons build a neighbour table; link-state reports build a graph; Dijkstra selects next hops; a per-hop priority queue orders transmissions; ACK, retry, deduplication and store-and-forward give reliability; AES-256-GCM and Ed25519 give confidentiality and authenticity; and an embedded web server exposes the node to a browser. The operating environment is an IP network segment on which infrastructure services have failed; where multicast is blocked a configured peer list is used (DIS-7). Dependencies are the Python runtime and the \texttt{cryptography} library.

\subsection{Functional Requirements}
\label{sec:fr}
The 68 functional requirements are listed in Table~\ref{tab:fr}. Each has a unique identifier, its text, and an acceptance criterion: the tests that must pass (and, for screens, a manual check). Each feature's input, processing and output are given in the feature's heading row.
\begingroup\footnotesize\renewcommand{\arraystretch}{1.0}
\setlength{\tabcolsep}{3pt}
\begin{longtable}{@{}p{1.25cm}p{6.3cm}p{1.65cm}p{6.95cm}@{}}
\caption[Functional requirements and their acceptance tests]{Functional requirements by feature (identifiers as in the SRS). Check: Auto = verified by automated tests, Auto+UI = automated tests plus a manual check of the screen, Partial = only partly implemented, Inspect = code inspection, Manual = manual check only; * marks a mapped test that failed once in repeated runs. The first mapped test is named and (+n) counts the others.}\\
\label{tab:fr}\\
\toprule \textbf{ID} & \textbf{Requirement (the system shall \ldots)} & \textbf{Check} & \textbf{Acceptance criterion: tests that must pass} \\ \midrule \endfirsthead
\toprule \textbf{ID} & \textbf{Requirement} & \textbf{Check} & \textbf{Acceptance criterion} \\ \midrule \endhead
\multicolumn{4}{@{}p{16.78cm}@{}}{\rule{0pt}{7pt}\textbf{3.4.1 Peer discovery.} \textit{Input:} HELLO beacons (role, public key) from neighbours; silence from a known neighbour. \textit{Processing:} Update the neighbour table; estimate link quality over a 20-HELLO window; prune after the timeout. \textit{Output:} Peer-up and peer-down events and link quality passed to routing.} \\ \midrule
DIS-1 & transmit a discovery beacon at a fixed, configurable interval. & Auto & \texttt{test\_\allowbreak{}hello\_\allowbreak{}is\_\allowbreak{}broadcast\_\allowbreak{}on\_\allowbreak{}start} \\
DIS-2 & keep a neighbour table with node id, address, role, public key, last-seen time and measured link quality. & Auto & \texttt{test\_\allowbreak{}update\_\allowbreak{}reports\_\allowbreak{}new\_\allowbreak{}neighbour\_\allowbreak{}once} (+3) \\
DIS-3 & add a previously unknown node to the table on receipt of a valid beacon. & Auto & \texttt{test\_\allowbreak{}neighbour\_\allowbreak{}added\_\allowbreak{}on\_\allowbreak{}hello\_\allowbreak{}received} (+1) \\
DIS-4 & declare a neighbour unreachable after a configurable period of silence. & Auto & \texttt{test\_\allowbreak{}neighbour\_\allowbreak{}marked\_\allowbreak{}dead\_\allowbreak{}after\_\allowbreak{}timeout} (+1) \\
DIS-5 & compute link quality as the proportion of expected beacons received over a sliding window. & Auto & \texttt{test\_\allowbreak{}perfect\_\allowbreak{}reception\_\allowbreak{}scores\_\allowbreak{}one} (+2) \\
DIS-6 & operate without manual configuration of peer addresses. & Auto & \texttt{test\_\allowbreak{}message\_\allowbreak{}to\_\allowbreak{}discovered\_\allowbreak{}neighbour\_\allowbreak{}delivers\_\allowbreak{}without\_\allowbreak{}manual\_\allowbreak{}next\_\allowbreak{}hop} (+1) \\
DIS-7 & support an optional configured seed list of peers where multicast is unavailable. & Auto & \texttt{test\_\allowbreak{}visible\_\allowbreak{}neighbours\_\allowbreak{}filters\_\allowbreak{}unlisted\_\allowbreak{}senders} (+1) \\
\multicolumn{4}{@{}p{16.78cm}@{}}{\rule{0pt}{7pt}\textbf{3.4.2 Message composition and transmission.} \textit{Input:} Text and destination, or the SOS control (sustained press). \textit{Processing:} Build a packet with a unique id and timestamp; encrypt and sign; queue. \textit{Output:} DATA packet sent; delivery state shown to the user.} \\ \midrule
MSG-1 & let a user compose and send a text message to a named node. & Auto+UI & \texttt{test\_\allowbreak{}send\_\allowbreak{}with\_\allowbreak{}ack\_\allowbreak{}delivers\_\allowbreak{}and\_\allowbreak{}gets\_\allowbreak{}acknowledged} (+3) \\
MSG-2 & let a user address a message to the rescue service rather than a specific node. & Auto & \texttt{test\_\allowbreak{}rescue\_\allowbreak{}dst\_\allowbreak{}routes\_\allowbreak{}to\_\allowbreak{}gateway\_\allowbreak{}across\_\allowbreak{}multiple\_\allowbreak{}hops} (+1) \\
MSG-3 & provide a distinct SOS control that requires a sustained press. & Auto+UI & \texttt{test\_\allowbreak{}sos\_\allowbreak{}reaches\_\allowbreak{}the\_\allowbreak{}rescue\_\allowbreak{}queue\_\allowbreak{}and\_\allowbreak{}the\_\allowbreak{}sender\_\allowbreak{}sees\_\allowbreak{}it\_\allowbreak{}delivered} \\
MSG-4 & assign a globally unique identifier to every message it originates. & Auto & \texttt{test\_\allowbreak{}msg\_\allowbreak{}id\_\allowbreak{}defaults\_\allowbreak{}to\_\allowbreak{}uuid4\_\allowbreak{}and\_\allowbreak{}round\_\allowbreak{}trips} (+1) \\
MSG-5 & record the creation time of every message it originates. & Auto & \texttt{test\_\allowbreak{}timestamp\_\allowbreak{}is\_\allowbreak{}uint64\_\allowbreak{}milliseconds} (+1) \\
MSG-6 & display the current delivery state of each originated message. & Auto+UI & \texttt{test\_\allowbreak{}store\_\allowbreak{}logs\_\allowbreak{}sent\_\allowbreak{}message\_\allowbreak{}and\_\allowbreak{}delivery\_\allowbreak{}lifecycle} (+1) \\
MSG-7 & display the number of intermediate nodes traversed by a delivered message. & Auto+UI & \texttt{test\_\allowbreak{}update\_\allowbreak{}message\_\allowbreak{}route\_\allowbreak{}records\_\allowbreak{}hops\_\allowbreak{}and\_\allowbreak{}path} (+1) \\
\multicolumn{4}{@{}p{16.78cm}@{}}{\rule{0pt}{7pt}\textbf{3.4.3 Multi-hop forwarding.} \textit{Input:} DATA or ACK packet from a neighbour. \textit{Processing:} Validate; drop duplicates; check hop allowance and loop; append own id; queue by priority. \textit{Output:} Packet forwarded, or dropped with a logged reason.} \\ \midrule
FWD-1 & forward packets not addressed to it towards the destination. & Auto & \texttt{test\_\allowbreak{}two\_\allowbreak{}hop\_\allowbreak{}relay\_\allowbreak{}a\_\allowbreak{}to\_\allowbreak{}b\_\allowbreak{}to\_\allowbreak{}c} (+1) \\
FWD-2 & decrement the remaining hop allowance of every packet it forwards. & Auto & \texttt{test\_\allowbreak{}on\_\allowbreak{}forward\_\allowbreak{}hook\_\allowbreak{}called\_\allowbreak{}with\_\allowbreak{}next\_\allowbreak{}hop} \\
FWD-3 & discard any packet whose hop allowance has reached zero. & Auto & \texttt{test\_\allowbreak{}drops\_\allowbreak{}when\_\allowbreak{}ttl\_\allowbreak{}expired} \\
FWD-4 & append its own identifier to the traversal path of every forwarded packet. & Auto & \texttt{test\_\allowbreak{}on\_\allowbreak{}forward\_\allowbreak{}hook\_\allowbreak{}called\_\allowbreak{}with\_\allowbreak{}next\_\allowbreak{}hop} (+1) \\
FWD-5 & discard any packet whose traversal path already contains its own identifier. & Auto & \texttt{test\_\allowbreak{}drops\_\allowbreak{}on\_\allowbreak{}loop\_\allowbreak{}detected} \\
FWD-6 & validate the structure of a packet before any forwarding decision. & Auto & \texttt{test\_\allowbreak{}bad\_\allowbreak{}magic\_\allowbreak{}rejected} (+3) \\
FWD-7 & be unable to read the payload of a packet it merely forwards. & Auto & \texttt{test\_\allowbreak{}relay\_\allowbreak{}never\_\allowbreak{}sees\_\allowbreak{}plaintext\_\allowbreak{}payload} \\
\multicolumn{4}{@{}p{16.78cm}@{}}{\rule{0pt}{7pt}\textbf{3.4.4 Adaptive routing.} \textit{Input:} Link-state reports; a destination. \textit{Processing:} Merge into the graph; Dijkstra over the weighted cost; resolve RESCUE\_\_ to the nearest gateway or rescue node. \textit{Output:} Least-cost next hop; routes recomputed on topology change.} \\ \midrule
RTE-1 & periodically transmit its own neighbour list to adjacent nodes. & Auto & \texttt{test\_\allowbreak{}send\_\allowbreak{}report\_\allowbreak{}reaches\_\allowbreak{}each\_\allowbreak{}neighbour} (+1) \\
RTE-2 & construct a graph of the network from received link-state reports. & Auto & \texttt{test\_\allowbreak{}receiving\_\allowbreak{}report\_\allowbreak{}merges\_\allowbreak{}reported\_\allowbreak{}edges\_\allowbreak{}into\_\allowbreak{}graph} (+1) \\
RTE-3 & compute least-cost paths using Dijkstra's algorithm over the weighted graph. & Auto & \texttt{test\_\allowbreak{}shortest\_\allowbreak{}path\_\allowbreak{}multi\_\allowbreak{}hop\_\allowbreak{}chain} (+1) \\
RTE-4 & derive link cost from a base hop penalty, measured link quality and recent delivery failures. & Auto & \texttt{test\_\allowbreak{}edge\_\allowbreak{}cost\_\allowbreak{}matches\_\allowbreak{}section\_\allowbreak{}8\_\allowbreak{}2\_\allowbreak{}formula} (+1) \\
RTE-5 & recompute routes whenever the topology changes. & Auto & \texttt{test\_\allowbreak{}record\_\allowbreak{}failure\_\allowbreak{}invalidates\_\allowbreak{}cached\_\allowbreak{}routes} (+2) \\
RTE-6 & treat any node advertising the gateway or rescue role as a valid terminal for rescue-service messages. & Auto & \texttt{test\_\allowbreak{}rescue\_\allowbreak{}dst\_\allowbreak{}resolves\_\allowbreak{}to\_\allowbreak{}nearest\_\allowbreak{}gateway} (+2) \\
RTE-7 & select the least-cost terminal when several are reachable. & Auto & \texttt{test\_\allowbreak{}rescue\_\allowbreak{}dst\_\allowbreak{}picks\_\allowbreak{}cheapest\_\allowbreak{}of\_\allowbreak{}multiple\_\allowbreak{}gateways} \\
RTE-8 & converge on a consistent topology view within a bounded interval after a change. & Auto & \texttt{test\_\allowbreak{}report\_\allowbreak{}reflooded\_\allowbreak{}reaches\_\allowbreak{}a\_\allowbreak{}node\_\allowbreak{}two\_\allowbreak{}hops\_\allowbreak{}away} (+1) \\
\multicolumn{4}{@{}p{16.78cm}@{}}{\rule{0pt}{7pt}\textbf{3.4.5 Priority handling.} \textit{Input:} Packets of different priority queued together. \textit{Processing:} Priority queue ordered by priority, then age, applied at every hop. \textit{Output:} SOS transmitted before routine traffic.} \\ \midrule
PRI-1 & support three priority levels: SOS, NORMAL, STATUS. & Auto & \texttt{test\_\allowbreak{}all\_\allowbreak{}three\_\allowbreak{}priorities\_\allowbreak{}ordered\_\allowbreak{}correctly} (+1) \\
PRI-2 & keep outgoing messages in a priority queue rather than a FIFO queue. & Auto & \texttt{test\_\allowbreak{}sos\_\allowbreak{}pops\_\allowbreak{}before\_\allowbreak{}normal\_\allowbreak{}regardless\_\allowbreak{}of\_\allowbreak{}insertion\_\allowbreak{}order} \\
PRI-3 & apply priority ordering at every hop, not only at the originator. & Auto & \texttt{test\_\allowbreak{}sos\_\allowbreak{}forwarded\_\allowbreak{}before\_\allowbreak{}already\_\allowbreak{}queued\_\allowbreak{}normal} (+1) \\
PRI-4 & order messages of equal priority by creation time, oldest first. & Auto & \texttt{test\_\allowbreak{}ties\_\allowbreak{}broken\_\allowbreak{}by\_\allowbreak{}oldest\_\allowbreak{}timestamp\_\allowbreak{}first} \\
PRI-5 & assign SOS priority to every message generated by the SOS control. & Auto+UI & \texttt{test\_\allowbreak{}sos\_\allowbreak{}reaches\_\allowbreak{}the\_\allowbreak{}rescue\_\allowbreak{}queue\_\allowbreak{}and\_\allowbreak{}the\_\allowbreak{}sender\_\allowbreak{}sees\_\allowbreak{}it\_\allowbreak{}delivered} \\
PRI-6 & display the priority level of each message in the interface. & Manual & \textit{Priority badges are shown on the screens (manual check); no automated UI test.} \\
\multicolumn{4}{@{}p{16.78cm}@{}}{\rule{0pt}{7pt}\textbf{3.4.6 Reliability.} \textit{Input:} Delivered or unacknowledged packets; duplicates; missing beacons. \textit{Processing:} ACK; timer; retry with a recomputed route; dedup set; failure detection. \textit{Output:} ACK returned; bounded retries; duplicates dropped; failed neighbours removed.} \\ \midrule
REL-1 & carry a hop allowance in every packet and decrement it at each hop. & Auto & \texttt{test\_\allowbreak{}drops\_\allowbreak{}when\_\allowbreak{}ttl\_\allowbreak{}expired} (+2) \\
REL-2 & remember recently observed message identifiers for a configurable period. & Auto & \texttt{test\_\allowbreak{}entry\_\allowbreak{}expires\_\allowbreak{}after\_\allowbreak{}window} (+1) \\
REL-3 & discard any packet whose identifier is already remembered. & Auto & \texttt{test\_\allowbreak{}duplicate\_\allowbreak{}delivery\_\allowbreak{}is\_\allowbreak{}dropped} (+2) \\
REL-4 & transmit an acknowledgement on successful delivery. & Auto & \texttt{test\_\allowbreak{}send\_\allowbreak{}with\_\allowbreak{}ack\_\allowbreak{}delivers\_\allowbreak{}and\_\allowbreak{}gets\_\allowbreak{}acknowledged} (+1) \\
REL-5 & retransmit an unacknowledged message after a configurable timeout. & Auto & \texttt{test\_\allowbreak{}timeout\_\allowbreak{}triggers\_\allowbreak{}retry\_\allowbreak{}and\_\allowbreak{}resets\_\allowbreak{}timer} (+1) \\
REL-6 & recompute the route before each retransmission. & Auto & \texttt{test\_\allowbreak{}record\_\allowbreak{}failure\_\allowbreak{}updates\_\allowbreak{}graph\_\allowbreak{}immediately\_\allowbreak{}not\_\allowbreak{}just\_\allowbreak{}on\_\allowbreak{}next\_\allowbreak{}report} (+1) \\
REL-7 & abandon delivery after a configurable number of failed attempts and mark the message accordingly. & Auto & \texttt{test\_\allowbreak{}exhausting\_\allowbreak{}max\_\allowbreak{}retries\_\allowbreak{}reports\_\allowbreak{}failed\_\allowbreak{}and\_\allowbreak{}stops\_\allowbreak{}tracking} (+1) \\
REL-8 & detect node failure through absence of discovery beacons and update the graph. & Auto & \texttt{test\_\allowbreak{}route\_\allowbreak{}forgotten\_\allowbreak{}when\_\allowbreak{}neighbour\_\allowbreak{}dies} (+2) \\
\multicolumn{4}{@{}p{16.78cm}@{}}{\rule{0pt}{7pt}\textbf{3.4.7 Store-and-forward.} \textit{Input:} A message with no route; a new route or neighbour. \textit{Processing:} Bounded retention by count and age; periodic flush; lowest priority evicted first. \textit{Output:} Message delivered when a path appears, or expired and logged.} \\ \midrule
SNF-1 & retain a message when route computation yields no path. & Auto & \texttt{test\_\allowbreak{}send\_\allowbreak{}without\_\allowbreak{}route\_\allowbreak{}queues\_\allowbreak{}to\_\allowbreak{}store\_\allowbreak{}forward} \\
SNF-2 & periodically re-attempt delivery of retained messages. & Auto & \texttt{test\_\allowbreak{}store\_\allowbreak{}forward\_\allowbreak{}flushes\_\allowbreak{}once\_\allowbreak{}a\_\allowbreak{}route\_\allowbreak{}appears} \\
SNF-3 & attempt delivery of retained messages immediately on detecting a new neighbour. & Auto & \texttt{test\_\allowbreak{}held\_\allowbreak{}message\_\allowbreak{}is\_\allowbreak{}sent\_\allowbreak{}the\_\allowbreak{}moment\_\allowbreak{}its\_\allowbreak{}neighbour\_\allowbreak{}appears} (+1) \\
SNF-4 & bound the retention queue by both message count and retention period. & Auto & \texttt{test\_\allowbreak{}enqueue\_\allowbreak{}over\_\allowbreak{}the\_\allowbreak{}bound\_\allowbreak{}evicts\_\allowbreak{}lowest\_\allowbreak{}priority} (+2) \\
SNF-5 & discard the lowest-priority retained messages first when the bound is reached. & Auto & \texttt{test\_\allowbreak{}enqueue\_\allowbreak{}over\_\allowbreak{}the\_\allowbreak{}bound\_\allowbreak{}evicts\_\allowbreak{}lowest\_\allowbreak{}priority} (+2) \\
SNF-6 & report the number of retained messages in the interface. & Manual & \textit{The held-message count is shown on the screens (manual check); no automated UI test.} \\
\multicolumn{4}{@{}p{16.78cm}@{}}{\rule{0pt}{7pt}\textbf{3.4.8 Message security.} \textit{Input:} Payload and packet identity fields. \textit{Processing:} AES-256-GCM encryption; Ed25519 signature; verification and decryption on receipt. \textit{Output:} Tampered or forged packets rejected.} \\ \midrule
SEC-1 & encrypt message payloads using authenticated symmetric encryption. & Auto & \texttt{test\_\allowbreak{}encrypt\_\allowbreak{}decrypt\_\allowbreak{}round\_\allowbreak{}trip} (+2) \\
SEC-2 & ensure intermediate nodes cannot recover the plaintext of messages they forward. & Auto & \texttt{test\_\allowbreak{}relay\_\allowbreak{}never\_\allowbreak{}sees\_\allowbreak{}plaintext\_\allowbreak{}payload} \\
SEC-3 & sign the packet header and traversal path with a digital signature scheme. & Auto & \texttt{test\_\allowbreak{}signed\_\allowbreak{}message\_\allowbreak{}verifies\_\allowbreak{}and\_\allowbreak{}delivers\_\allowbreak{}one\_\allowbreak{}hop} (+2) \\
SEC-4 & verify the signature of every received packet before delivery. & Auto & \texttt{test\_\allowbreak{}forged\_\allowbreak{}signature\_\allowbreak{}is\_\allowbreak{}rejected\_\allowbreak{}not\_\allowbreak{}delivered} (+1) \\
SEC-5 & discard any packet failing signature verification or authentication. & Auto & \texttt{test\_\allowbreak{}forged\_\allowbreak{}signature\_\allowbreak{}is\_\allowbreak{}rejected\_\allowbreak{}not\_\allowbreak{}delivered} (+2) \\
SEC-6 & distribute node public keys through discovery beacons. & Auto & \texttt{test\_\allowbreak{}hello\_\allowbreak{}includes\_\allowbreak{}public\_\allowbreak{}key\_\allowbreak{}when\_\allowbreak{}configured} (+1) \\
SEC-7 & use established cryptographic libraries and not implement primitives directly. & Inspect & \textit{Code inspection: crypto.py only calls the cryptography library (AESGCM, Ed25519).} \\
\multicolumn{4}{@{}p{16.78cm}@{}}{\rule{0pt}{7pt}\textbf{3.4.9 Logging and history.} \textit{Input:} Every message and network event. \textit{Processing:} Write to the node's SQLite database. \textit{Output:} History and delivery statistics available by query.} \\ \midrule
LOG-1 & maintain a persistent local record of messages sent, received and forwarded. & Auto & \texttt{test\_\allowbreak{}store\_\allowbreak{}logs\_\allowbreak{}forward\_\allowbreak{}on\_\allowbreak{}middle\_\allowbreak{}node} (+1) \\
LOG-2 & record status, creation time, resolution time, hop count and path of each message. & Auto & \texttt{test\_\allowbreak{}record\_\allowbreak{}message\_\allowbreak{}stores\_\allowbreak{}path\_\allowbreak{}as\_\allowbreak{}comma\_\allowbreak{}separated} (+2) \\
LOG-3 & record network events: node arrival and failure, route change, duplicate suppression, hop-allowance expiry, retransmission, retention. & Auto & \texttt{test\_\allowbreak{}store\_\allowbreak{}logs\_\allowbreak{}peer\_\allowbreak{}up\_\allowbreak{}and\_\allowbreak{}down} (+2) \\
LOG-4 & store plaintext content only where the node is the originator or destination. & Inspect & \texttt{test\_\allowbreak{}store\_\allowbreak{}logs\_\allowbreak{}forward\_\allowbreak{}on\_\allowbreak{}middle\_\allowbreak{}node} \textit{Forwarded rows have body NULL (Table~\ref{tab:samples}); asserted for the relay's log, otherwise by inspection.} \\
LOG-5 & let the user view and filter the local message history. & Auto+UI & \texttt{test\_\allowbreak{}history\_\allowbreak{}filters\_\allowbreak{}by\_\allowbreak{}priority\_\allowbreak{}and\_\allowbreak{}direction} (+2) \\
LOG-6 & support extraction of delivery statistics from the stored records. & Auto & \texttt{test\_\allowbreak{}list\_\allowbreak{}and\_\allowbreak{}count\_\allowbreak{}messages\_\allowbreak{}filter\_\allowbreak{}by\_\allowbreak{}direction\_\allowbreak{}and\_\allowbreak{}status} (+1) \\
\multicolumn{4}{@{}p{16.78cm}@{}}{\rule{0pt}{7pt}\textbf{3.4.10 Rescue dashboard.} \textit{Input:} Node state and events. \textit{Processing:} Compute topology, SOS queue and statistics from the graph and the logs. \textit{Output:} Live graph; SOS queue by priority then age; statistics.} \\ \midrule
DSH-1 & render the mesh topology as a graph of nodes and active links. & Auto+UI & \texttt{test\_\allowbreak{}snapshot\_\allowbreak{}describes\_\allowbreak{}a\_\allowbreak{}converged\_\allowbreak{}mesh} \\
DSH-2 & update the topology view as nodes join and fail. & Auto+UI & \texttt{test\_\allowbreak{}killing\_\allowbreak{}a\_\allowbreak{}node\_\allowbreak{}reroutes\_\allowbreak{}and\_\allowbreak{}reviving\_\allowbreak{}restores\_\allowbreak{}it} \\
DSH-3 & display outstanding emergency messages ordered by priority, then age. & Auto+UI & \texttt{test\_\allowbreak{}sos\_\allowbreak{}reaches\_\allowbreak{}the\_\allowbreak{}rescue\_\allowbreak{}queue\_\allowbreak{}and\_\allowbreak{}the\_\allowbreak{}sender\_\allowbreak{}sees\_\allowbreak{}it\_\allowbreak{}delivered} \\
DSH-4 & display the originating node and traversal path of each received message. & Auto+UI & \texttt{test\_\allowbreak{}sos\_\allowbreak{}reaches\_\allowbreak{}the\_\allowbreak{}rescue\_\allowbreak{}queue\_\allowbreak{}and\_\allowbreak{}the\_\allowbreak{}sender\_\allowbreak{}sees\_\allowbreak{}it\_\allowbreak{}delivered} (+1) \\
DSH-5 & display aggregate statistics: delivery ratio, mean hop count, nodes online. & Auto+UI & \texttt{test\_\allowbreak{}snapshot\_\allowbreak{}describes\_\allowbreak{}a\_\allowbreak{}converged\_\allowbreak{}mesh} \\
DSH-6 & receive updates without the coordinator refreshing the view. & Auto+UI & \texttt{test\_\allowbreak{}websocket\_\allowbreak{}streams\_\allowbreak{}a\_\allowbreak{}hello\_\allowbreak{}snapshot\_\allowbreak{}then\_\allowbreak{}events} \\
\bottomrule
\end{longtable}
\endgroup


\subsection{Non-Functional Requirements}
\label{sec:nfr}
\begin{itemize}
\item \textbf{Performance:} at least 90\,\% delivery on a forty-node network with 20\,\% packet loss; convergence of the topology view within 15\,s; delivery resumed within 10\,s of a node failure; an SOS measurably faster than routine traffic under load.
\item \textbf{Scalability:} at least fifty nodes without degradation beyond these targets and without configuration changes to existing nodes.
\item \textbf{Availability and reliability:} no single point of failure; correct operation under packet loss, node failure and partition; no message discarded while retention capacity remains.
\item \textbf{Security and privacy:} payloads encrypted; relays cannot read them; tampering detected; a forged SOS rejected; every packet field validated; private keys never transmitted; plaintext stored only on the originator and destination.
\item \textbf{Usability and accessibility:} the emergency action is one prominent control requiring a sustained press; no end-user configuration; plain-language status text and a state change visible within 2\,s.
\item \textbf{Maintainability, portability and resources:} independent modules for each concern; the engine is independent of the link technology, so another transport replaces the current one unchanged; one process on one event-loop thread with an embedded database and bounded queues (measured in Chapter~\ref{ch:perf}).
\end{itemize}

\subsection{Hardware and Software Requirements}
\textbf{Hardware:} laptops (Windows~11 used for development and testing) on a Wi-Fi network or phone hotspot; no server. \textbf{Software:} Python~3.10 with \texttt{asyncio}; \texttt{cryptography}~42 (AES-256-GCM, Ed25519); SQLite in WAL mode; FastAPI, uvicorn, pydantic and \texttt{websockets} for the API; React~18, Vite~6, Tailwind~4 and d3-force for the web application (with a manifest and service worker so it installs as a PWA); pytest, pytest-asyncio and httpx for testing; Git and GitHub for version control.

\subsection{Constraints and Assumptions}
\textbf{Constraints:} one semester (phases from 17~August to 2~November 2026); three members; no radio hardware, so the radio link is emulated by a neighbour-restricted LAN; scale about 50 nodes; no machine learning and no native mobile application. \textbf{Assumptions:} the software is installed before an incident; every device is within range of another, directly or through intermediaries; batteries last the incident; at least one gateway or an arriving rescue node exists for messages to leave the mesh; node identifiers are unique and assigned in advance.

\subsection{Use Cases and User Stories}
\textbf{UC1:} as a trapped user I hold the SOS control so that rescuers learn I need help; the system sends the highest-priority message and shows its live status. \textbf{UC2:} as a trapped user with no route I want my message kept; it is retained, flushed when a path appears, and the screen says it is held. \textbf{UC3:} as a relay user I want to help without reading others' messages; the node forwards automatically and the payload is opaque to it. \textbf{UC4:} as a gateway operator I want to be an exit point; the node advertises the gateway role. \textbf{UC5:} as a rescue coordinator I want the SOS queue ordered by priority then age, with route and age, and the live topology. \textbf{UC6:} as an administrator I want to restrict who hears whom through a topology file. The use case diagram is Figure~\ref{fig:usecase}.
